\section*{Chapter \ref{ch:s2:systematic-credit-and-bond-etf-arbitrage} --- Systematic Credit and Bond-ETF Arbitrage}

\begin{solution}{exo:s2:systematic-credit-and-bond-etf-arbitrage:1}
About 300 basis points: the price is $0.9496$ of NAV and the NAV $1.0215$ of true value, so the price is $0.9496 \times 1.0215 = 0.970$ of true
value. The two parts multiply rather than add, which is why $215 + 300$ exceeds 504.
\end{solution}

\begin{solution}{exo:s2:systematic-credit-and-bond-etf-arbitrage:2}
Buying at $0.97$ of true value and selling the bonds at true value earns $1/0.97 - 1 = 3.09\%$ of the redeemed value; less the cost of 1.5\% that is
159 basis points. The chapter's best day, 163, adds the price noise.
\end{solution}

\begin{solution}{exo:s2:systematic-credit-and-bond-etf-arbitrage:3}
$1.50\% \times 6.5 = 9.75\%$, close to the simulated $9.8\%$; carry over three weeks and the spread of durations around 6.5 account for the rest.
\end{solution}

\begin{solution}{exo:s2:systematic-credit-and-bond-etf-arbitrage:4}
A bond's excess return is mostly its spread duration times its spread change, and spreads rise with duration. A signal compared across all bonds
would buy long bonds for their wide spreads and load on market credit risk, which is a beta, not a factor. Comparing against peers of similar
duration and rating, and balancing the long and short sides in spread duration, isolates the characteristic.
\end{solution}

\begin{solution}{exo:s2:systematic-credit-and-bond-etf-arbitrage:5}
Haddad, Moreira and Muir found that investors tried to raise cash by selling their safer and more liquid holdings. The selling fell on the ETFs of
Treasuries, municipal and investment-grade bonds, which investors hold as liquid reserves, more than on high-yield ETFs.
\end{solution}

\begin{solution}{exo:s2:systematic-credit-and-bond-etf-arbitrage:6}
The index holds thousands of bonds, many of which rarely trade; delivering or receiving all of them would be costly or impossible. Koont, Ma, Pastor
and Zeng found that bond ETFs choose baskets with cash and a subset of bonds, the more so for illiquid ones, and adjust them to correct imbalances
between creations and redemptions while still allowing arbitrage.
\end{solution}

\begin{solution}{exo:s2:systematic-credit-and-bond-etf-arbitrage:7}
With bonds trading on 5\% of days, the worst discount doubles to 1,050 basis points, still fourteen trading days into the stress, and the stale part
rises from 215 to 838 basis points: the NAV lags true value by about nineteen days on average, longer than the stress, so it misses most of the
widening. The part that is selling pressure is unchanged, since the price follows true value, but a discount against NAV now mostly measures stale
prices.
\end{solution}

\begin{solution}{exo:s2:systematic-credit-and-bond-etf-arbitrage:8}
The NAV is computed from stale prices; in the chapter's stress 215 of the 504 basis points is NAV that has not caught up with falling bonds, and it
closes by the NAV falling, not the price rising. Only the part against true value, 300 basis points, is a return, and only to someone who can redeem
and sell the bonds, at a cost of 150 basis points in the stress.
\end{solution}

\begin{solution}{pb:s2:systematic-credit-and-bond-etf-arbitrage:1}
\begin{enumerate}
\item A characteristic of bonds, such as spread against peers, spread momentum, duration or size, that sorts them into portfolios with persistently
different duration-adjusted excess returns.
\item Size, low-risk, value and momentum portfolios earned economically meaningful and statistically significant alphas; low correlations made a
combination more efficient; the results held after costs and in liquid bonds.
\item 300 bonds over eight years, durations from one to twelve; spreads move with a market factor, a persistent issuer trend, a reverting gap and
noise; expected excess returns rise with the square root of duration. Value is the spread above the fitted spread-duration line less its
three-month change; momentum is the six-month tightening; low risk buys short bonds against long ones at equal risk.
\item Value 0.46, momentum 0.14, low risk 0.95, combined 0.94; correlations $-0.09$ (value, momentum), $-0.03$ (value, low risk) and $-0.02$
(momentum, low risk).
\item Create-redeem arbitrage: delivering a basket for new shares when shares trade rich, or shares for the basket when they trade cheap, keeping the
gap less costs. The discount: price below published NAV.
\item Bond ETFs choose baskets with cash and a subset of the index, especially when bonds are illiquid, and adjust them to correct imbalances; basket
inclusion helps liquidity in general but hurts it in large imbalances such as the COVID-19 crisis.
\item The mean of the bonds' last traded values; each bond trades on a quarter of days.
\item A trade of a whole list of bonds at a single negotiated price, priced against the ETF and the bonds' quotes and hedgeable with the ETF.
\item Six years in, market spreads widen 150 basis points over fifteen days and half of it returns over forty; the ETF price falls a further 3\% below
true value at the depth and recovers over forty days. The basket loses 9.8\%.
\item At the worst, $-504$ basis points against NAV: stale NAV 215 above true value, price 300 below it, multiplied.
\item On 28 days, all in the stress, by up to 163 basis points of the redeemed value after a cost of 1.5\%; never outside the stress, where the
discount averages 2.7 basis points with a standard deviation of 21.6.
\item Liquid bond ETFs traded at large discounts to NAV in March 2020, more for Treasuries, municipal and investment-grade bonds than high yield, as
investors sold safer assets for cash; the disruptions reversed after the Fed's announcements of 23 March and 9 April.
\item Sharpe ratios of 0.95 (low risk), 0.46 (value), 0.14 (momentum) and 0.94 combined; in the stress the ETF traded 5.0\% below NAV and 3.0\% below
true value, and the authorised participant could profit on 28 days by up to 1.63\%.
\item True value is right, and no one sees it because most bonds did not trade: the NAV is stale and the price is pushed by sellers.
\item The reverting gap drives part of each six-month spread change, so momentum partly buys bonds whose gap is about to revert, against value.
\item From the ETF price, the prices of bonds that did trade and their co-movement with those that did not, and dealer quotes, with a model of how
far selling pressure moves the ETF.
\item Against an estimate of true value from traded prices, not the stale NAV; with basket rules as they were and costs of selling bonds in stress.
\item The ETF premium-discount arbitrage.
\item Chapter~18's negative basis is bonds below their CDS because funding is scarce; this chapter's discount is an ETF below its bonds because cash
is scarce; both widen in the same crises and pay whoever has balance sheet left.
\item How much the bonds' true value has moved since they last traded, plus how much sellers of the ETF will pay for immediacy.
\end{enumerate}
\end{solution}

\begin{solution}{iq:s2:systematic-credit-and-bond-etf-arbitrage:1}
When the shares trade above the basket's value, an authorised participant buys the basket, delivers it for new shares and sells them; below it, the
participant buys shares, redeems them for the basket and sells the basket. Both trades push the price back towards the basket's value, within the
participant's costs.
\end{solution}

\begin{solution}{iq:s2:systematic-credit-and-bond-etf-arbitrage:2}
The spread against a fitted curve of peers of similar rating, duration and sector, possibly against a model of default risk; remove the part that is
recent widening, balance the sides in spread duration, check it is not just default risk, and price it with realistic bond costs.
\end{solution}

\begin{solution}{iq:s2:systematic-credit-and-bond-etf-arbitrage:3}
Against the ETF and the bonds' quotes: hedge the list at once with the ETF, estimate how long each bond takes to work out and what it costs, add the
risk of the list diverging from the ETF over that time, and charge more for bonds the ETF does not hold.
\end{solution}

\begin{solution}{iq:s2:systematic-credit-and-bond-etf-arbitrage:4}
The ETF falls below its bonds and the book is marked at a loss; the bonds cannot be sold at their marks, redemptions deliver baskets that may not
match the book, and funding tightens. The hedge works only for whoever can hold until the discount closes.
\end{solution}

\begin{solution}{iq:s2:systematic-credit-and-bond-etf-arbitrage:5}
Start from each bond's last trade and dealer quotes; move it with the observed moves of liquid proxies (the ETF, CDS, the bonds of the same issuer
and similar bonds that did trade) weighted by its duration and rating; flag its age and confidence, and test the estimate against the next trade.
\end{solution}

\begin{solution}{iq:s2:systematic-credit-and-bond-etf-arbitrage:6}
The last trade was $k$ days ago with probability $p(1-p)^k$, so the NAV reflects true value $k$ days old, and $E[k] = \sum_k k\, p(1-p)^k =
(1-p)/p$. After a steady move of $m$ a day that has lasted longer than most bonds' gaps between trades, the stale gap is $m\,E[k] = m(1-p)/p$:
three days' move for $p = 1/4$, nineteen for $p = 1/20$.
\end{solution}
